Three extended falsification checks probe the construct validity of the AI proxy, the direction of the AI-TFP relationship, and the sensitivity of the result to the small complete-case sample.

First, replacing the composite AI index with a digital-infrastructure-only sub-index (internet use, mobile subscriptions, and secure servers) yields a coefficient of 0.036 (p-value 0.028, n = 2053). An innovation-only sub-index (resident and non-resident patent applications and IP receipts) yields 0.018 (p-value below 0.001, n = 1832). Both sub-indices are positively associated with TFP, indicating that the headline association is not isolated to AI-specific activity.

Second, a reverse-causality check regresses log AI adoption on one-period-lagged log TFP. The estimated coefficient is -1.558 (p-value 0.367, n = 1452).

Third, restricting the baseline sample to countries with at least 8 of 15 years of non-missing AI data (n = 1433 observations) yields an AI coefficient of -0.010 (p-value below 0.001).

For context, only 163 of the 193 countries in the cleaned panel report AI data in at least one year, contributing 1898 country-year observations with non-missing AI data, against 997 country-years without it. Table~\ref{tab:sample-selection} compares observable characteristics across these two groups.